\begin{tcolorbox}[slot, enhanced, colback=white, colframe=ink, boxrule=1.6pt, arc=4mm,
  left=12pt, right=12pt, top=6pt, bottom=7pt,
  title={%
    {\fontsize{52}{52}\selectfont\bfseries\itshape i\kern0.04em}%
    {\fontsize{40}{44}\selectfont\bfseries\scshape n the Beginning was the Math}\hfill
    {\fontsize{22}{26}\selectfont\itshape $i^2=-1$\,; engineers call it $j$, for $i$ is betrothed.}},
  colbacktitle=ink, coltitle=white, toptitle=4pt, bottomtitle=4pt, lefttitle=10pt, righttitle=14pt]
\mathbandsize\setdisplay
% #1 width, #2 content; the three columns are named so the slides can rearrange them
\newcommand{\mathcol}[2][0.327\linewidth]{\begin{minipage}[t]{#1}\vspace{0pt}#2\end{minipage}}
\newcommand{\mathA}{%
  \begin{eqgroup}[Euler: the most important line in the degree]
  \[ e^{j\theta}=\cos\theta+j\sin\theta \]
  \end{eqgroup}
  \begin{eqgroup}[A sinusoid is a complex number]
  \[ A\cos(\omega t+\varphi)=\mathrm{Re}\{\underbrace{Ae^{j\varphi}}_{I+jQ}\,e^{j\omega t}\} \]
  \end{eqgroup}
  \begin{eqgroup}[Taylor: why linearisation works]
  \[ f(x_0+\Delta x)\approx f(x_0)+f'(x_0)\,\Delta x \]
  \end{eqgroup}
  \begin{eqgroup}[Decibels: logs turn products into sums]
  \[ G_{\mathrm{dB}}=10\log_{10}\frac{P_2}{P_1}=20\log_{10}\frac{V_2}{V_1} \]
  \end{eqgroup}
  % the slide has room the poster hasn't
  \slidesonly{%
  \begin{eqgroup}[Hilbert: signals are vectors, energy is length]
  \[ \langle x,y\rangle=\int x(t)\,y^*(t)\,dt \]
  \[ \|x\|^2=\langle x,x\rangle \]
  \[ |\langle x,y\rangle|\le\|x\|\,\|y\| \quad\text{\small\color{muted}(why the matched filter wins)} \]
  \end{eqgroup}}}
\newcommand{\mathB}{%
  \begin{eqgroup}[Linear systems and their modes]
  \[ A\mathbf x=\mathbf b \]
  \[ A\mathbf v=\lambda\mathbf v \]
  \[ \det(A-\lambda I)=0 \]
  \end{eqgroup}
  \begin{eqgroup}[Transforms turn calculus into algebra]
  \[ X(\omega)=\int x(t)\,e^{-j\omega t}\,dt \]
  \[ \mathcal L\{\dot y\}=sY(s)-y(0^-) \]
  \[ X(z)=\sum_n x[n]\,z^{-n} \qquad z=e^{sT_s} \]
  \[ \tfrac{d}{dt}\to j\omega:\ \ Z_L=j\omega L \quad Z_C=1/j\omega C \]
  \end{eqgroup}
  \begin{eqgroup}[Orthogonality: every coefficient is a projection]
  \[ c_k=\frac{\langle x,\varphi_k\rangle}{\langle\varphi_k,\varphi_k\rangle} \]
  \end{eqgroup}}
\newcommand{\mathC}{%
  \begin{eqgroup}[Divergence and Stokes]
  \[ \oint_S\mathbf F\cdot d\mathbf A=\int_V\nabla\cdot\mathbf F\,dV \]
  \[ \oint_C\mathbf F\cdot d\mathbf l=\int_S(\nabla\times\mathbf F)\cdot d\mathbf A \]
  \end{eqgroup}
  \begin{eqgroup}[Chance: Bayes, sums and the bell]
  \[ P(A\,|\,B)=\frac{P(B\,|\,A)\,P(A)}{P(B)} \]
  \[ \mathrm{Var}(X+Y)=\mathrm{Var}\,X+\mathrm{Var}\,Y \quad\text{(indep.)} \]
  \[ \bar x\ \sim\ \mathcal N\!\big(\mu,\ \sigma^2/n\big) \quad\text{(CLT)} \]
  \end{eqgroup}
  \begin{eqgroup}[Parseval and Dirac]
  \[ \int|x(t)|^2dt=\frac{1}{2\pi}\int|X(\omega)|^2d\omega \]
  \[ \int x(\tau)\,\delta(t-\tau)\,d\tau=x(t) \]
  \end{eqgroup}}
\newcommand{\maths}{\mathcol{\mathA}\hfill\mathcol{\mathB}\hfill\mathcol{\mathC}}
\newcommand{\secondorder}{%
  \begin{eqgroup}[Second order: time and frequency]
  \[ \ddot y+2\zeta\omega_0\dot y+\omega_0^2y=\omega_0^2u \]
  \end{eqgroup}}
% on the slides the Bode plot sits beside the step response instead of under it, and both widen to fill the row
% (the legend then stands in the gap between them)
\ifslides\pgfplotsset{bode at/.style={at={(step.south east)}, anchor=south west, xshift=44mm}, resp width/.style={width=136mm}, step height/.style={height=36mm}, bode height/.style={height=34mm},
  resp legend/.style={legend style={at={(1.075,0.62)}, anchor=west, row sep=-1pt}}}
\else\pgfplotsset{bode at/.style={at={(step.south west)}, anchor=north west, yshift=-9mm}, resp width/.style={width=118mm}, step height/.style={height=29mm}, bode height/.style={height=27mm},
  resp legend/.style={legend style={at={(1.03,-0.15)}, anchor=west, row sep=-1pt}}}\fi
\newcommand{\responses}{%
  \begin{tikzpicture}
  \begin{axis}[posterplot, name=step, scale only axis, resp width, step height, xmin=0, xmax=10.5, ymin=0, ymax=1.6,
     xlabel={$\omega_0 t$}, ylabel={step}, every axis y label/.style={at={(axis description cs:0,1)}, anchor=south west, xshift=2mm},
     xtick=\empty, ytick={1}, samples=140, domain=0:10,
     legend style={draw=none, fill=none, font=\small}, resp legend, legend cell align=left]
    \addplot[accent, smooth] {1-exp(-0.2*x)*(cos(deg(0.9798*x))+0.2041*sin(deg(0.9798*x)))};
    \addlegendentry{$\zeta=0.2$}
    \addplot[accentlt, smooth] {1-exp(-0.7*x)*(cos(deg(0.7141*x))+0.9802*sin(deg(0.7141*x)))};
    \addlegendentry{$\zeta=0.7$}
    \addplot[good, smooth] {1-(1+x)*exp(-x)};
    \addlegendentry{$\zeta=1$}
    \addplot[trap, smooth] {1+(-3.732*exp(-0.268*x)+0.268*exp(-3.732*x))/3.464};
    \addlegendentry{$\zeta=2$}
    \draw[muted, dashed, thin] (axis cs:0,1) -- (axis cs:10.5,1);
  \end{axis}
  \begin{semilogxaxis}[posterplot, name=bode, bode at, scale only axis, axis lines=left, resp width, bode height, xmin=0.1, xmax=1000, ymin=-126, ymax=12,
     xlabel={$\omega/\omega_0$}, ylabel={$|H|$ [dB]},
     every axis y label/.style={at={(axis description cs:0,0)}, anchor=south west, xshift=3mm, yshift=1mm},
     every axis x label/.style={at={(axis description cs:1,0)}, anchor=north east},
     xtick={1,10,100}, xticklabels={$1$,$10$,$100$}, ytick=\empty, samples=180, domain=0.1:1000]
    \addplot[accent, smooth] {-10*log10((1-x^2)^2+(0.4*x)^2)};
    \addplot[accentlt, smooth] {-10*log10((1-x^2)^2+(1.4*x)^2)};
    \addplot[good, smooth] {-10*log10((1-x^2)^2+(2*x)^2)};
    \addplot[trap, smooth] {-10*log10((1-x^2)^2+(4*x)^2)};
    \node[muted, font=\small, anchor=south west] at (axis cs:30,-45) {$-40$ dB/dec};
  \end{semilogxaxis}
  \posteronly{\node[rotate=-90, anchor=south, text=muted, font=\small, align=center]
    at ($(step.above north east)!0.5!(bode.below south east)+(30mm,0)$) {RLC, servo, suspension, PLL:\\ one equation.};}
  \end{tikzpicture}}
% slides: the second order closes the first column; the plots sit under the other two
\ifslides
\mathcol{\mathA\vspace{4mm}
  \secondorder\vspace{2mm}
  {\centering\small\color{muted}RLC, servo, suspension, PLL: one equation.\par}}\hfill
\begin{minipage}[t]{0.665\linewidth}\vspace{0pt}
  \mathcol[0.4917\linewidth]{\mathB}\hfill\mathcol[0.4917\linewidth]{\mathC}\par
  \vspace{6mm}{\centering\responses\par}
\end{minipage}\par
\else
\begin{minipage}[t]{0.66\linewidth}
\maths
\end{minipage}\hfill
\begin{minipage}[t]{0.325\linewidth}\vspace{0pt}\centering
  \secondorder\vspace{1mm}
  \responses
\end{minipage}
\fi

\vspace{6pt}
\trap{Behold, every signal hath two faces, time and frequency: cross over unto the side where the problem is easy, and ye shall find rest. \textbf{Prove all things:} check thy units, thy limits and thy signs; \textbf{hold fast that which is good.}}
\end{tcolorbox}
